\documentclass[11pt]{article}
\usepackage[utf8]{inputenc}
\usepackage[T1]{fontenc}
\usepackage{amsmath,amssymb}
\usepackage{booktabs}
\usepackage{hyperref}
\usepackage{url}
\usepackage{xcolor}
\usepackage{geometry}
\usepackage{enumitem}
\usepackage{natbib}

\geometry{a4paper, margin=1in}

\title{When Embedding-Based Methods Cannot Distinguish Hallucinations from Truth: A Human-Confabulated Benchmark}

\author{Javier Mar\'{i}n \\
\texttt{javier@jmarin.info}}

\date{}

\begin{document}
\maketitle

\begin{abstract}
Large language models generate outputs that sound correct but are not---a phenomenon generally called hallucination, confabulation, or unfaithful generation.
Detecting these outputs is both a research challenge and a regulatory requirement: the EU AI Act mandates accuracy measurement for high-risk AI systems.
We report a practical finding encountered while developing geometric detection methods for production RAG systems: existing hallucination benchmarks, predominantly built by prompting LLMs to generate false content, did not provide a clear signal for differentiating grounded from ungrounded—hallucinations—responses in embedding space.
Our hypothesis is that the benchmarks carry generation artifacts---distributional signatures of the instruction-conditioned generation mode---that detectors learn instead of learning to identify hallucination itself.
To test this hypothesis and address our practical need, we construct an alternative benchmark: 212 question--response pairs across nine domains where a human without domain expertise writes plausible-sounding but entirely false responses.
We analyze five representative examples and show that intuitive human confabulation strategies---redefinition of terms, inversion of mechanisms, invention of entities, anchoring in false regulatory histories, and displacement across semantic categories---each exploit a specific property of how autoregressive language models represent text.
These properties encode linguistic form but not referential truth, establishing a theoretical ceiling for embedding-based detection.
We release the dataset and building protocol to support the evaluation of hallucination detection methods against this ceiling.
\end{abstract}

% ============================================================
\section{Measuring Accuracy in High-Risk AI Systems}
\label{sec:problem}

Article~15 of the EU Artificial Intelligence Act \citep{euaiact2024} requires that providers of high-risk AI systems implement metrics to measure output accuracy.\footnote{EU AI Act, Article 15 (page 17): ``High-risk AI systems should perform consistently throughout their lifecycle and meet an appropriate level of
accuracy, robustness and cybersecurity, in light of their intended purpose and in accordance with the generally acknowledged state of the art.'' \url{https://eur-lex.europa.eu/legal-content/EN/TXT/PDF/?uri=OJ:L_202401689}}
For systems built on LLMs---particularly retrieval-augmented generation (RAG) pipelines in regulated industries like healthcare, legal, and financial services---this creates a clear engineering requirement: the system must detect when its output is wrong.

We use the term \emph{confabulation} specifically.
\citet{berlyne1972confabulation} introduced the clinical definition: the production of false statements without the intention to deceive---what he called ``honest lying.''
In neuropsychology, confabulation refers to the production of false information without the intention to deceive---the confabulator does not know that the content is false and fills gaps in knowledge with plausible-sounding material \citep{moscovitch1995confabulation, dallabarba1993different}.
\citet{farquhar2024detecting} adopted this term for LLM outputs, distinguishing confabulations (arbitrary and incorrect generations) from deliberate deception.
Our construction protocol operationalizes the same process: the author lacks domain knowledge and fills that gap with plausible-sounding content, producing text that is structurally coherent and referentially empty.

We approached confabulations geometrically. Modern sentence transformers are trained via contrastive objectives that organize text on the unit hypersphere $\mathbb{S}^{d-1}$ \citep{reimers2019sentence, wang2020understanding}, where angular distance serves as the natural metric \citep{mardia2000directional}.
If confabulated responses---outputs that sound correct but misrepresent the facts---occupy different angular positions than grounded responses, the geometry provides a detection signal.
In prior work \citep{marin2025sgi}, we formalized this approach as the Semantic Grounding Index (SGI), which measures the ratio of angular distances from a response to the query versus the retrieved context.

SGI works when context is available. But many production deployments involve direct prompting, conversational interfaces, or scenarios where the retrieval step is opaque.
Without a context anchor, detection requires calibration: a reference set of verified grounded and ungrounded responses from which the detector learns what each looks like in a given domain.

\section{Hallucination Benchmark Landscape}
\label{sec:existing}

The hallucination evaluation literature has grown rapidly, with benchmarks addressing different tasks, domains, and definitions of hallucination \citep{ji2023survey, huang2025survey}. We review the construction methods most relevant to our geometric detection setting.

\subsection{LLM-generated benchmarks.}
The largest hallucination benchmarks generate false content by prompting LLMs.
HaluEval \citep{li2023halueval} uses ChatGPT in a two-step sampling-and-filtering pipeline to produce 35{,}000 samples. HaluBench \citep{ravi2024lynx} combines existing datasets with synthetic perturbations for 15{,}000 RAG-focused triplets. HalluLens \citep{bang2025hallulens} introduces dynamic test generation from Wikipedia to prevent data contamination.

In these benchmarks, the false content is generated by an autoregressive model conditioned on an instruction to confabulate. In an autoregressive transformer, every token is generated from $p(x_t \mid x_{<t})$, where the preceding context $x_{<t}$ includes all prior tokens — including any instruction tokens. When a prompt such as ``generate a hallucinated answer'' is prepended to the query, those instruction tokens enter the attention computation: their key–value pairs influence every subsequent generated token through every layer \citep{vaswani2017attention}. The output distribution is therefore mathematically different from what the same model would produce without the instruction, because the conditioning context is different. Whether this distributional difference introduces detectable signatures — artifacts that a detector could learn independently of the hallucination content — is an empirical question.

Two lines of evidence suggest it does. \citet{gudibande2023false} documented an analogous confound in model imitation. Models trained on ChatGPT outputs learn to reproduce style — fluency, structure, register — without improving factual accuracy, proving that distributional properties of instruction-conditioned text are separable from factual content. Independently, \citet{bang2025hallulens} showed that existing hallucination benchmarks, including HaluEval, conflate hallucination with factuality and do not adequately evaluate extrinsic hallucinations — content inconsistent with the model's training data rather than with a provided input context. If the benchmarks themselves do not cleanly isolate hallucination, detectors trained or evaluated on them may learn to exploit confounds in the benchmark construction rather than hallucination itself.

\subsection{Natural output collection.}
FactScore \citep{min2023factscore} evaluates factual precision by decomposing natural LLM outputs into atomic facts verified against Wikipedia.
RAGTruth \citep{niu2024ragtruth} collects 18{,}000 naturally generated RAG responses with word-level human annotations.
FaithBench \citep{bao2025faithbench} selects summarization outputs where automated detectors disagree, then applies expert annotation.
These approaches avoid the generation-artifact confound by evaluating what models produce naturally, but they cannot control for specific confabulation types, require expensive annotation, and---critically---do not provide the paired grounded/ungrounded examples needed for geometric calibration.

\subsection{Misconception targeting.}
TruthfulQA \citep{lin2022truthfulqa} includes 817 questions targeting common misconceptions. The false answers are not novel confabulations but \emph{widely shared errors}---beliefs that many humans hold and that LLMs reproduce from training data.
For domain-specific detection in finance, medicine, or law, we needed confabulations that resemble what models actually produce in production: confident, domain-appropriate text that is wrong in domain-specific ways, not in ways that reflect popular misconceptions.

\subsection{Sampling-based detection.}
An alternative approach bypasses benchmarks entirely.
SelfCheckGPT \citep{manakul2023selfcheckgpt} and semantic entropy methods \citep{kuhn2023semantic, farquhar2024detecting} detect confabulation by sampling multiple responses and measuring consistency. True knowledge produces stable outputs and confabulated content varies across samples.
These methods require 5--20 generations per query, making them impractical for production-scale deployment where a large amount of queries require real-time evaluation.

\section{Building a Human-Confabulated Benchmark}
\label{sec:construction}

\subsection{Protocol}

We built the dataset using a method motivated by practical constraints rather than theoretical design. The benchmark has four properties: 
\begin{enumerate}
    \item Paired grounded and ungrounded responses to the same questions, 
    \item Covers specialized domains where confabulation causes real harm, 
    \item The confabulated content resembles real production confabulations rather than benchmark artifacts, and 
    \item The confabulated content is novel---not drawn from training data or popular belief.
\end{enumerate}

We created pairs of QA samples with a factually correct answer generated by a capable LLM (Claude Sonnet 4.5) and manually verified against authoritative sources. Confabulated responses---hallucinations--- are written by the authors following a single instruction: to write a response that would sound convincing to someone who does not know the subject, without looking up the answer, inventing every factual claim. The response had to sound like something a domain non-expert would accept without suspicion---the kind of confident, authoritative output that characterizes production LLM confabulations in regulated industries.

\subsection{Domain Coverage and Statistics}

The dataset includes 212 pairs across nine domains (Table~\ref{tab:domains}), selected to span knowledge types encountered in production deployments: regulatory knowledge (finance, law), clinical procedures (medicine), declarative fact (science, geography, history), technical specification (Python, TypeScript), and general knowledge.

\begin{table}[t]
\centering
\small
\begin{tabular}{lrl}
\toprule
\textbf{Domain} & \textbf{$n$} & \textbf{Knowledge type} \\
\midrule
Python coding    & 47 & Technical specification \\
Finance          & 40 & Regulatory / procedural \\
Medical          & 40 & Clinical / declarative \\
Science          & 21 & Declarative fact \\
TypeScript coding & 18 & Technical specification \\
History          & 14 & Declarative fact \\
Law              & 11 & Regulatory / procedural \\
General knowledge & 11 & Mixed \\
Geography        & 10 & Declarative fact \\
\midrule
\textbf{Total} & \textbf{212} & \\
\bottomrule
\end{tabular}
\caption{Domain distribution of the benchmark.}
\label{tab:domains}
\end{table}

Response lengths were not explicitly controlled; however, the natural tendency to elaborate when confabulating produced comparable distributions (grounded: $\mu = 55.6$ words, median $= 55.0$; confabulated: $\mu = 57.3$, median $= 57.0$), preventing length from serving as a trivial detection signal.

\section{Dataset Analysis}
\label{sec:analysis}

The confabulations were written intuitively.
But analyzing them reveals structural patterns that explain why they approach a detection ceiling for embedding-based methods. We examine an example per knowledge type---5 in total---.

\subsection{Redefinition within the register}

\begin{quote}
\textbf{Q:} What does P/E ratio stand for in stock market analysis? \\[5pt]
\textbf{Grounded:} ``P/E ratio stands for Price-to-Earnings ratio. It is calculated by dividing a company's current stock price by its earnings per share.'' \\[5pt]
\textbf{Confabulated:} ``The P/E ratio stands for Price-to-Exit ratio and is commonly used by Private Equity firms to valuate their investment performance.''
\end{quote}

This confabulation redefines the acronym while remaining entirely within the financial register.
``P/E ratio,'' ``Private Equity,'' ``firms,'' ``investment,'' and ``performance'' all co-occur naturally in financial text.
An embedding model trained on distributional co-occurrence \citep{harris1954distributional, firth1957synopsis} will assign similar representations to both responses, because the distributional context is almost identical.
The difference is referential: Price-to-Earnings exists; Price-to-Exit (in this context) does not. Referential correspondence is not a distributional property---it requires access to the world beyond the text.

This connects to the distinction \citet{bender2020climbing} drew between form--- any observable realization of language--- and meaning---the relation between the form and
something external to language---. In this case, the confabulation preserves the form of a financial definition while replacing the meaning entirely.

\subsection{Mechanism inversion}

\begin{quote}
\textbf{Q:} What is photosynthesis? \\[5pt]
\textbf{Grounded:} ``Photosynthesis is the process by which green plants [...] absorb carbon dioxide from the air and water from the soil, producing glucose and oxygen.'' \\[5pt]
\textbf{Confabulated:} ``Photosynthesis is the chemical process by which plants absorb oxygen from the atmosphere and convert it into nitrogen compounds that fertilize the surrounding soil.''
\end{quote}

In this example, the confabulation inverts the direction of the mechanism---oxygen absorbed rather than produced; nitrogen compounds rather than glucose---while preserving every local transition. Autoregressive models predict the next token from preceding tokens: $p(x_t | x_{<t})$ \citep{vaswani2017attention}. Each transition in the confabulation---``plants'' $\to$ ``absorb,'' ``absorb'' $\to$ ``oxygen,'' ``convert'' $\to$ ``into'' $\to$ ``nitrogen'' $\to$ ``compounds''---follows plausible local distributional patterns.
The error is global: the mechanism as a whole does not exist. But the training objective operates on local transitions and provides no explicit signal for global factual coherence.

\subsection{Entity invention through composition}

\begin{quote}
\textbf{Q:} Which organ produces insulin? \\[5pt]
\textbf{Grounded:} ``The pancreas produces insulin. This organ contains specialized cells called beta cells, located in clusters known as the islets of Langerhans.'' \\[5pt]
\textbf{Confabulated:} ``Insulin is produced exclusively by the brain, specifically in the hypothalamus region which monitors blood glucose levels continuously through specialized neural receptors.''
\end{quote}

The confabulation replaces the correct organ (pancreas) with one that has a genuine metabolic role (hypothalamus) and invents a supporting mechanism (``specialized neural receptors''). Transformer representations compose meanings from constituent tokens \citep{mikolov2013linguistic}: ``hypothalamus'' + ``monitors'' + ``blood glucose'' is a distributionally valid composition because each pairwise relationship exists in medical text. The composed meaning---hypothalamus producing insulin---does not correspond to reality, but compositional validity does not require referential accuracy \citep{bender2020climbing}.

\subsection{Reinterpretation through polysemy}

\begin{quote}
\textbf{Q:} What is \textit{habeas corpus}? \\[5pt]
\textbf{Grounded:} ``\textit{habeas corpus} is a legal principle that protects individuals from unlawful detention by requiring the government to justify holding a person in custody before a court.'' \\[5pt]
\textbf{Confabulated:} ``\textit{habeas corpus }is a Latin legal term meaning \emph{to have the body of evidence}, referring to the prosecution's obligation to physically present all material evidence in court.''
\end{quote}

The Latin expression``\textit{habeas corpus}'' literally means ``you shall have the body.'' The confabulation reinterprets ``body'' from its legal sense (the detained person) to an alternative sense (body of evidence)---exploiting polysemy.
Contextual embedding models resolve polysemy through attention to surrounding tokens \citep{devlin2019bert, vaswani2017attention}: when the context consistently supports the wrong sense (``evidence,'' ``prosecution,'' ``physically present''), the model resolves the ambiguity accordingly. The resulting representation is internally consistent but grounded in the wrong interpretation.

\subsection{Template-filling in structured domains}

\begin{quote}
\textbf{Q:} What does the \texttt{\_\_init\_\_} method do in a Python class? \\[5pt]
\textbf{Grounded:} ``The \texttt{\_\_init\_\_} method is a special method in Python classes that serves as the constructor. It is automatically called when a new instance of the class is created.'' \\[5pt]
\textbf{Confabulated:} ``The \texttt{\_\_init\_\_} method is a private system function that initializes Python's garbage collector at the start of every script execution. It runs automatically before the main module loads and allocates the base memory pool.''
\end{quote}

The confabulation preserves the structural template of a technical specification---a method name, its classification (``private system function''), its trigger (``runs automatically''), and its effect (``allocates... memory pool'')---while replacing every specific claim. Domains with rigid discourse templates — coding documentation, legal statutes, clinical guidelines — are particularly vulnerable to distributional confabulation. In these domains, the template itself (a method name, its classification, its trigger condition, its effect) carries the main distributional signal. It determines which tokens are likely to co-occur, in what order, and with what syntactic relationships. The specific content that fills the template slots — where factual truth resides — contributes comparatively little to the resulting embedding, because the distributional neighborhood is already largely determined by the template structure. A confabulation that preserves the template while replacing every slot value occupies nearly the same region of embedding space as the truthful response.


\section{Discussion}
\label{sec:discussion}

\subsection{Why can't embedding-based methods distinguish some hallucinations from truth?}
The five examples analyzed in Section~\ref{sec:analysis} illustrate that each confabulation strategy preserves a different subset of the distributional properties encoded by transformer-based embedding models — lexical co-occurrence statistics, local transition probabilities, compositional phrase structure, contextual polysemy resolution, and domain-specific discourse templates — while systematically violating referential correspondence with the external world. Embedding-based methods cannot reliably distinguish these confabulations from the truth because truth is not a distributional property. Detection requires referential information---correspondence between text and the world---which lies outside what distributional representations capture \citep{harris1954distributional, bender2020climbing}.


This does not say that embedding-based detection is useless. Many real-world confabulations do introduce distributional anomalies — unfamiliar vocabulary, atypical structures, inconsistent register — and methods such as SelfCheckGPT \citep{manakul2023selfcheckgpt}, FActScore \citep{min2023factscore}, and semantic entropy \citep{farquhar2024detecting} exploit these signals effectively when these anomalies are present. But the ceiling defines the boundary beyond which distributional methods cannot reach, and progress on the hard cases requires approaches that go beyond distributional similarity, retrieval-augmented verification, mechanistic interpretability, or methods that access model internals not available in the embedding space.

\subsection{The generation artifact hypothesis}
Our empirical findings are consistent with the hypothesis that LLM-generated benchmarks carry detectable artifacts from the instruction-conditioned generation mode: on such benchmarks, the embeddings of grounded and ungrounded---hallucinated---responses occupy geometrically distinct regions, and detection methods achieve high accuracy. On our human confabulations, this geometric difference collapses. Grounded and confabulated responses are nearly indistinguishable in embedding space. Human confabulations remove the generation-mode confound entirely. Whether this difference is robust across embedding models and detection methods is an empirical question that we hope the released dataset will help the community investigate.

\subsection{Limitations}
The dataset was produced by a single confabulator, whose stylistic patterns may introduce person-specific artifacts analogous to the generation artifacts we critique; a multi-confabulator extension would test generalization. Also, the dataset is English only. The example analysis in Section~\ref{sec:analysis} was performed post-hoc---the structural patterns were identified after the dataset was built, not before.
We report that the construction was intuitive, and the analysis is our attempt to understand why the intuitions produced effective confabulations. Future work should quantify the distributional similarity between grounded and confabulated responses across multiple embedding models and compare detection performance against LLM-generated benchmarks under controlled conditions. Our preliminary results across seven embedding models with different architectures and a different number of parameters suggest that the pattern holds (results will be reported separately in different papers).


\section{Conclusion}

We built a hallucination benchmark for detecting confabulation in general and domain-specific LLM deployments, as required for accuracy measurement under EU AI Act Article~15. This dataset exploits the structural properties of how LLMs represent text: distributional co-occurrence, local transition patterns, compositional representation, contextual polysemy resolution, and domain-specific templates.
Each property encodes form without encoding truth, explaining why these confabulations approach the theoretical ceiling for embedding-based detection. We release the 212-pair dataset and construction protocol at \url{https://github.com/Javihaus/cert-confabulation-benchmark} to contribute to the evaluation of detection methods at and above this ceiling.


\bibliographystyle{plainnat}

\begin{thebibliography}{28}

\bibitem[Bang et al.(2025)]{bang2025hallulens}
Bang, Y., et al.
\newblock HalluLens: LLM Hallucination Benchmark.
\newblock In \emph{Proceedings of the 63rd Annual Meeting of the Association for Computational Linguistics (Volume 1: Long Papers), pages 24128--24156}, 2025.

\bibitem[Bao et al.(2025)]{bao2025faithbench}
Bao, Y., et al.
\newblock FaithBench: A Diverse Hallucination Benchmark for Summarization by Modern LLMs.
\newblock In \emph{Proceedings of the 2025 Conference of the Nations of the Americas Chapter of the Association for Computational Linguistics: Human Language Technologies (Volume 2: Short Papers)}, pages 448--461, 2025.

\bibitem[Bender and Koller(2020)]{bender2020climbing}
Bender, E.~M. and Koller, A.
\newblock Climbing towards NLU: On Meaning, Form, and Understanding in the Age of Data.
\newblock In \emph{Proceedings of the 58th Annual Meeting of the Association for Computational Linguistics. Association for Computational Linguistics}, pages 5185--5198,Online, 2020.

\bibitem[Berlyne(1972)]{berlyne1972confabulation}
Berlyne, N.
\newblock Confabulation.
\newblock \emph{British Journal of Psychiatry}, 120(554):31--39, 1972.

\bibitem[Dalla~Barba(1993)]{dallabarba1993different}
Dalla~Barba, G.
\newblock Different patterns of confabulation.
\newblock \emph{Cortex}, 29(4):567--581, 1993.

\bibitem[Devlin et al.(2019)]{devlin2019bert}
Devlin, J., Chang, M.-W., Lee, K., and Toutanova, K.
\newblock BERT: Pre-training of deep bidirectional transformers for language understanding.
\newblock In \emph{Proceedings of the 2019 conference of the North American chapter of the association for computational linguistics: human language technologies, volume 1 (long and short papers)}, pages 4171--4186, 2019.

\bibitem[European Parliament(2024)]{euaiact2024}
European Parliament and Council.
\newblock Regulation (EU) 2024/1689 of 13 June 2024 laying down harmonized rules on artificial intelligence (AI Act).
\newblock \emph{Official Journal of the European Union}, L series, 2024.

\bibitem[Farquhar et al.(2024)]{farquhar2024detecting}
Farquhar, S., Kossen, J., Kuhn, L., and Gal, Y.
\newblock Detecting hallucinations in large language models using semantic entropy.
\newblock \emph{Nature}, 630, 625--630, 2024.

\bibitem[Firth(1957)]{firth1957synopsis}
Firth, J.~R.
\newblock A synopsis of linguistic theory 1930--1955.
\newblock In \emph{Studies in Linguistic Analysis}, pages 1--32. Blackwell, 1957.

\bibitem[Gudibande et al.(2023)]{gudibande2023false}
Gudibande, A., et al.
\newblock The False Promise of Imitating Proprietary LLMs.
\newblock \emph{arXiv preprint arXiv:2305.15717}, 2023.

\bibitem[Harris(1954)]{harris1954distributional}
Harris, Z.~S.
\newblock Distributional structure.
\newblock \emph{Word}, 10(2-3):146--162, 1954.

\bibitem[Huang et al.(2025)]{huang2025survey}
Huang, L., et al.
\newblock A Survey on Hallucination in Large Language Models: Principles, Taxonomy, Challenges, and Open Questions.
\newblock In \emph{ACM Transactions on Information Systems}, 43(2), 1-55., 2025.

\bibitem[Ji, Z et al.(2023)]{ji2023survey}
Ji, Z., et al.
\newblock Survey of hallucination in natural language generation
\newblock In \emph{ACM computing surveys}, 55(12), pages 1-38, 2023.

\bibitem[Kuhn et al.(2023)]{kuhn2023semantic}
Kuhn, L., Gal, Y., and Farquhar, S.
\newblock Semantic uncertainty: Linguistic invariances for uncertainty estimation in natural language generation.
\newblock \emph{arXiv preprint arXiv:2302.09664}, 2023.

\bibitem[Li et al.(2023)]{li2023halueval}
Li, J., Cheng, X., Zhao, W.~X., Nie, J.-Y., and Wen, J.-R.
\newblock HaluEval: A Large-Scale Hallucination Evaluation Benchmark for Large Language Models.
\newblock In \emph{Proceedings of the 2023 Conference on Empirical Methods in Natural Language Processing}, pages 6449–6464, Singapore, 2023.

\bibitem[Lin et al.(2022)]{lin2022truthfulqa}
Lin, S., Hilton, J., and Evans, O.
\newblock TruthfulQA: Measuring How Models Mimic Human Falsehoods.
\newblock In \emph{Proceedings of the 60th annual meeting of the association for computational linguistics (volume 1: long papers)}, pages 3214-3252, 2022.

\bibitem[Manakul et al.(2023)]{manakul2023selfcheckgpt}
Manakul, P., Liusie, A., and Gales, M.~J.~F.
\newblock SelfCheckGPT: Zero-Resource Black-Box Hallucination Detection for Generative Large Language Models.
\newblock In \emph{Proceedings of the 2023 conference on empirical methods in natural language processing}, pages 9004-9017, 2023.

\bibitem[Mardia and Jupp(2000)]{mardia2000directional}
Mardia, K.~V. and Jupp, P.~E.
\newblock \emph{Directional Statistics}.
\newblock John Wiley \& Sons, 2000.

\bibitem[Mar\'{i}n(2025)]{marin2025sgi}
Mar\'{i}n, J.
\newblock Semantic Grounding Index: Geometric Bounds on Context Engagement in RAG Systems.
\newblock \emph{arXiv preprint arXiv:2512.13771}, 2025.

\bibitem[Maynez et al.(2020)]{maynez2020faithfulness}
Maynez, J., Narayan, S., Bohnet, B., and McDonald, R.
\newblock On Faithfulness and Factuality in Abstractive Summarization.
\newblock In \emph{Proceedings of the 58th Annual Meeting of the Association for Computational Linguistics}, pages 1906--1919, Online, 2020.

\bibitem[Mikolov et al.(2013)]{mikolov2013linguistic}
Mikolov, T., Yih, W., and Zweig, G.
\newblock Linguistic regularities in continuous space word representations.
\newblock In \emph{Proceedings of NAACL}, pages 746--751, 2013.

\bibitem[Min et al.(2023)]{min2023factscore}
Min, S., et al.
\newblock FActScore: Fine-grained Atomic Evaluation of Factual Precision in Long Form Text Generation.
\newblock In \emph{Proceedings of the 2023 Conference on Empirical Methods in Natural Language Processing.}, pages 12076-12100, 2023.

\bibitem[Moscovitch(1995)]{moscovitch1995confabulation}
Moscovitch, M.
\newblock Confabulation.
\newblock In Schacter, D.~L., editor, \emph{Memory Distortions}, pages 226--251. Harvard University Press, 1995.

\bibitem[Niu et al.(2024)]{niu2024ragtruth}
Niu, C., et al.
\newblock RAGTruth: A Hallucination Corpus for Developing Trustworthy Retrieval-Augmented Language Models.
\newblock In \emph{Proceedings of the 62nd Annual Meeting of the Association for Computational Linguistics (Volume 1: Long Papers)}, pages 10862-10878, 2024.

\bibitem[Ravi et al.(2024)]{ravi2024lynx}
Ravi, S., et al.
\newblock Lynx: An Open Source Hallucination Evaluation Model.
\newblock \emph{arXiv preprint arXiv:2407.08488}, 2024.

\bibitem[Reimers and Gurevych(2019)]{reimers2019sentence}
Reimers, N. and Gurevych, I.
\newblock Sentence-BERT: Sentence Embeddings using Siamese BERT-Networks.
\newblock In \emph{Proceedings of the 2019 conference on empirical methods in natural language processing and the 9th international joint conference on natural language processing (EMNLP-IJCNLP)}, pages 3982-3992, 2019.

\bibitem[Vaswani et al.(2017)]{vaswani2017attention}
Vaswani, A., et al.
\newblock Attention Is All You Need.
\newblock In \emph{Advances in Neural Information Processing Systems}, 2017.

\bibitem[Wang and Isola(2020)]{wang2020understanding}
Wang, T. and Isola, P.
\newblock Understanding Contrastive Representation Learning through Alignment and Uniformity on the Hypersphere.
\newblock \emph{arXiv preprint arxiv:2005.10242v10}, 2020.

\end{thebibliography}

\end{document}
